%% ma: L11-B14-0015
%% lop: 11
%% chuong: 4
%% bai: 14
%% dang: 11.14.4
%% loai: TN4
%% muc: TH
%% dap_an: "A"
%% nguon: "tai-lieu-hoc-tap-mon-toan-lop-11-theo-chuong-trinh-giao-duc-pho-thong-2018.pdf (Lớp Toán Thầy Duy, Phạm Lê Duy, THPT Chu Văn An), Chương 4 Bài 5, Phần 1, Câu 10 (tr.528)"
%% kien_thuc: "Lớp 11 Bài 14 (phép chiếu song song)"
%% trang_thai: chinh-thuc
%% ghi_chu: "Đã duyệt 03/11/2026."
\begin{ex}
Cho hình chóp $S.ABCD$ đáy là hình bình hành tâm $O$. Trên cạnh $SB$, $SD$ lần lượt lấy điểm $M,N$ sao cho $SM=2MB$ và $SN=\dfrac13SD$. Hình chiếu của $M,N$ qua phép chiếu song song lên mặt phẳng chiếu $(ABCD)$ theo phương của đường thẳng $SO$ lần lượt là $P,Q$. Tính tỉ số $\dfrac{OP}{OQ}$.
\choice{\True $2$.}{$\dfrac12$.}{$\dfrac34$.}{$\dfrac23$.}
\loigiai{\sloppy Theo phương $SO$: $S\mapsto O$, $B\mapsto B$, $D\mapsto D$. Phép chiếu bảo toàn tỉ số trên cùng đường thẳng. $\dfrac{OP}{OB}=\dfrac{SM}{SB}=\dfrac23$ và $\dfrac{OQ}{OD}=\dfrac{SN}{SD}=\dfrac13$. Vì $OB=OD$ nên $\dfrac{OP}{OQ}=\dfrac{2/3}{1/3}=2$.}
\end{ex}
